\documentclass[book.tex]{subfiles}

\begin{document}

\chapter{The Heat Equation}

\label{chap:heat_equation}
\noindent\rule{11cm}{0.4pt} \\
\textbf{Prerequisites:} \autoref{chap:pde}, \autoref{chap:simple_physics} \\
\textbf{Difficulty Level:} ** \\
\noindent\rule{11cm}{0.4pt}
\vspace{5mm}
The heat equation is a fundamental partial differential equation that models systems with dissipating heat.

Let $u(x,y,z,t)$ be a function of spatial coordinates and time. The heat equation is given by
\begin{equation}
    \frac{\partial u}{\partial t} = \alpha \left ( \frac{\partial^2 u}{\partial x^2} + \frac{\partial^2 u}{\partial y^2} + \frac{\partial^2 u}{\partial z^2} \right )
\end{equation}

\begin{equation}
    \frac{\partial u}{\partial t} = \alpha \nabla^2 u
\end{equation}

The $\nabla^2$ operator is a general "averaging" operator. So the heat equation takes a point and moves it to the average of its neighbors over time.

We're going to take some time to study solutions to the heat equation in some depth. While the transmission of heat is indeed an interesting problem, the heat equation is also the simplex example of a physical partial differential equation that we will encounter in this book so it serves as a neat test bed to learn the core principles of solving physical equations.

\subsection{Linearity}
One of the most important properties of the heat equation is that its solutions satisfy a linearity property. That is, assume that $u_1$ and $u_2$ are both solutions. Then we know that 
\begin{align}
    \frac{\partial u_1}{\partial t} &= \alpha \nabla^2 u_1 \\
    \frac{\partial u_2}{\partial t} &= \alpha \nabla^2 u_2 
\end{align}
By the fact that partial derivatives and $\nabla^2$ are both linear operators, we find that a linear combination of $u_1$ and $u_2$ is also a solution
\begin{equation}
    \frac{\partial (c_1 u_1 + c_2 u_2)}{\partial t} = \alpha \nabla^2(c_1 u_1 + c_2 u_2)
\end{equation}
This observation provides a powerful tool for solving the heat equation in greater generality, since if we can find some solutions, we can find more solutions by linear combinations.

\subsection{Separation of Variables}
The idea of separation of variables is simple. We're going to posit that there exists a solution to the heat equation which has the following special structure.

\begin{equation}
    u(x, t) = X(x)T(t)
\end{equation}

That is, the solution is "separable". Note that if we have a family of separable solutions of the form $X_i(x)T_i(t)$, we can construct more general solutions by linear combinations
\begin{equation}
    u(x, t) = \sum_i X_i(x)T_i(t)
\end{equation}
Note that $u$ need not be separable now! So the separable solutions form a sort of "basis" which we can use to solve for the general case of this equation.

Let's now plug in separable equations into the heat equation

\begin{align}
    \frac{\partial u}{\partial t} &= \alpha \nabla^2 u \\
    X(x) T'(t) &= \alpha T(t) \nabla^2 X \\
    \frac{T'(t)}{T(t)} &= \alpha \frac{(\nabla^2 X)(x)}{X(x)}
\end{align}
This last set of equations has an interesting form. Note that the left hand side is a function of $t$ while the right hand side is a function of $x$. As $x$ and $t$ independently vary, this equality must hold. The only way this can be the case is if both sides equal constants. So we can factor this equation into two separate equations
\begin{align}
    \frac{T'(t)}{T(t)} &= c \\
    \alpha \frac{(\nabla^2 X)(x)}{X(x)} &= c
\end{align}
Multiplying through we recover the two ordinary differential equations
\begin{align}
    T'(t) &= cT(t) \\
    \alpha (\nabla^2 X)(x) &= cX(x) 
\end{align}
For simplicity, let's work with these equations in 1D. Then our equations reduce to
\begin{align}
    T'(t) &= cT(t) \\
    X''(x) &= \frac{c}{\alpha}X(x) \label{1d_X_eqn}
\end{align}
The solution to the equation for $T$ should recognizably be the exponential function
\begin{equation}
    T(t) = d e^{ct}
    \label{1d_X_soln}
\end{equation}
The solution for the order two differential equation is a little trickier though. We can start from a guess that the solution perhaps should look like $e^{\lambda t}$ since we have that $X''$ and $X$ are multiple of one another. Let's try substituting in $X = e^{\lambda t}$ into our equation and seeing what happens
\begin{align}
    \lambda^2 e^{\lambda t} &= \frac{c}{\alpha} e^{\lambda t} \\
    \lambda^2 &= \frac{c}{\alpha}
\end{align}
Note that cancelling the $e^{\lambda t}$ term is allowed since this term is never 0. We now see that equality holds when
\begin{equation}
    \lambda = \pm \sqrt{\frac{c}{\alpha}}
\end{equation}
Note that Equation \ref{1d_X_eqn} satisfies the property (like the heat equation itself) that its solutions space is linear. It follows that we have a family of equations
\begin{equation}
    X(x) = c_1 e^{\sqrt{\frac{c}{\alpha}}t} + c_2 e^{-\sqrt{\frac{c}{\alpha}}t}
\end{equation}
In fact, every solution to Equation \ref{1d_X_eqn} is of this form, but we won't prove that fact here. We now have a family of separable solutions for the heat equation in 1D
\begin{equation}
    u(x, t) = \left (c_1 e^{\sqrt{\frac{c}{\alpha}}t} + c_2 e^{-\sqrt{\frac{c}{\alpha}}t} \right ) \left (d e^{ct} \right)
    \label{heat_1d_general}
\end{equation}

\section{Boundary Conditions}
While we've identified a general class of solutions to the heat equation, in practice we'll have some additional information about the physical system that we're trying to model. This type of information in general is termed \textit{boundary conditions.}.

In this section, we'll work through a simple 1D example of boundary conditions to see how we can turn the generic solution form in Equation \ref{heat_1d_general} into solutions that meet our boundary conditions. In particular, let's suppose that we have a 1D rod whose endpoints are constrained to have a fixed temperature of $0$ for all time

\begin{align}
u(0, t) &= 0 \qquad \textrm{for all } t > 0 \\
u(\ell, t) &= 0 \qquad \textrm{for all } t > 0
\end{align}

We'll also assume that we know the initial temperature of the rod everywhere.

\begin{equation}
    u(x, 0) = f(x) \qquad \textrm{for all } 0 < x < \ell
\end{equation}

Let's now impose our boundary conditions onto our general solutions. Let's start with the condition that the boundary temperature must be $0$ for all time. Note that under our separation of variables setting, we can enforce boundary conditions by requiring that $X(0) = 0$ and $X(\ell) = 0$ for a separable solution $X(x)T(t)$. Note that a more general solution will have the form
\begin{equation}
    u(x, t) = \sum_i X_i(x)T_i(t)
\end{equation}
If we enforce $X_i(0) = 0$ and $X_i(\ell) = 0$, we will have that $u$ satisfies the endpoint conditions for boundary temperature. %Now how about the condition for the known initial temperature? 

%We're going to take a bit of a leap now. We know that linear combinations of equations of the form in Equation \ref{heat_1d_general} are also solutions. 

\section{Exercises}
\begin{enumerate}
    \item Given the differential equation $T'(t) = cT(t)$ verify that 
    \begin{equation}
        T(t) = d \exp{ct}
    \end{equation}
    is a solution where $d$ is an arbitrary constant. Prove that all solutions must be of this form.
    \item Given the differential equation
    \begin{equation}
        X''(x) = \frac{c}{\alpha} X
    \end{equation}
    Show that if $f_1$ and $f_2$ are solutions to this equation, so too is $c_1 f_1 + c_2 f_2$
\end{enumerate}

\end{document}